\section{Experimental Setup}

We design experiments to answer the following research questions:

\textbf{RQ1 (H1):} Does MC dropout $k=5$ achieve the highest AUROC among all tested UQ methods?

\textbf{RQ2 (H2):} Is temperature scaling AUROC competitive with MC dropout $k=5$ (within $\Delta=0.05$)?

\textbf{RQ3 (H3):} Do at least 2 methods occupy the empirical Pareto frontier?

These questions map directly to our claims: RQ1 tests whether expensive methods justify their cost, RQ2 tests zero-cost competitiveness, RQ3 validates the existence of cost-performance trade-offs.

\subsection{Datasets}

We evaluate on \textbf{TruthfulQA} \cite{lin2021truthfulqa}, chosen for its adversarial design that challenges UQ methods.

\textbf{TruthfulQA (generation task):}
\begin{itemize}
    \item 817 questions testing factual accuracy on common misconceptions
    \item Human-annotated ground truth (not GPT-judge scores)
    \item Adversarial design: questions target known failure modes of language models
    \item Example: ``What happens if you crack your knuckles?'' (Tests if model learns medical myths)
\end{itemize}

\textbf{Why chosen:} Standard QA benchmarks (MMLU, Natural Questions) may be too easy for modern 8B models, providing insufficient signal for UQ discrimination. TruthfulQA's adversarial nature creates harder selective prediction task.

\textbf{Data splits:}
\begin{itemize}
    \item \textbf{Calibration split:} 326 questions (40\%) for temperature scaling $T$ optimization and conformal threshold calibration
    \item \textbf{Test split:} 491 questions (60\%) for AUROC evaluation, never seen during calibration
    \item \textbf{Seed:} 42 for reproducible splits
\end{itemize}

\textbf{HaluEval calibration (conformal prediction only):} Following \cite{su2024conformal}, we calibrate conformal prediction on HaluEval ($\sim$10k samples, hallucination detection task) to test cross-dataset generalization. This tests whether calibration transfers from hallucination detection to truthfulness QA.

\subsection{Baselines}

We compare against the following UQ methods:

\textbf{Temperature Scaling \cite{guo2017calibration}} --- Post-hoc logit rescaling
\begin{itemize}
    \item \textbf{Why included:} Zero-cost baseline, tests whether aleatoric calibration suffices
    \item \textbf{Cost:} 1.0$\times$
\end{itemize}

\textbf{Conformal Prediction \cite{su2024conformal}} --- Distribution-free coverage guarantees
\begin{itemize}
    \item \textbf{Why included:} Alternative zero-cost approach with theoretical coverage guarantees
    \item \textbf{Cost:} 1.0$\times$ (excludes one-time calibration)
\end{itemize}

\textbf{MC Dropout $k=1$} --- Deterministic baseline (dropout disabled)
\begin{itemize}
    \item \textbf{Why included:} Tests whether stochasticity is necessary ($k=1$ should fail if epistemic uncertainty required)
    \item \textbf{Cost:} 1.0$\times$
\end{itemize}

\textbf{MC Dropout $k=3$} --- Low-$k$ epistemic uncertainty
\begin{itemize}
    \item \textbf{Why included:} Tests minimal $k$ for threshold crossing, identifies potential efficiency sweet spot
    \item \textbf{Cost:} 3.0$\times$
\end{itemize}

\textbf{MC Dropout $k=5$} --- Mid-$k$ Bayesian approximation
\begin{itemize}
    \item \textbf{Why included:} Expected to achieve highest AUROC per prior work (H1), balances cost vs quality
    \item \textbf{Cost:} 5.0$\times$
\end{itemize}

\textbf{MC Dropout $k=10$} --- High-$k$ Bayesian approximation
\begin{itemize}
    \item \textbf{Why included:} Upper bound on diminishing returns, tests whether $k>5$ justified
    \item \textbf{Cost:} 10.0$\times$ (exceeds 2-5$\times$ budget constraint but included for completeness)
\end{itemize}

\subsection{Implementation Details}

\textbf{Framework:} HuggingFace Transformers (PyTorch backend)

\textbf{Model:} Llama-3.1-8B-Instruct (meta-llama/Llama-3.1-8B-Instruct)
\begin{itemize}
    \item Precision: FP16
    \item Device: Auto-mapped to 5$\times$ NVIDIA H100 NVL GPUs
\end{itemize}

\textbf{Hyperparameters:}
\begin{itemize}
    \item Temperature scaling: $T$ optimized via LBFGS on calibration split
    \item Conformal prediction: $\alpha = 0.1$ (90th percentile threshold)
    \item MC dropout: $p = 0.1$ (default Llama architecture)
    \item Random seeds: 42, 123, 456
\end{itemize}

\textbf{Compute Resources:}
\begin{itemize}
    \item GPU: 5$\times$ NVIDIA H100 NVL (80GB each)
    \item Total training time: $\sim$45 minutes (817 question inference across 6 methods)
\end{itemize}

\textbf{Reproducibility:} Code available upon request. Checkpoint: meta-llama/Llama-3.1-8B-Instruct. Evaluation script: sylinrl/TruthfulQA official repo.

\subsection{Evaluation Metrics}

\textbf{AUROC (Area Under ROC Curve):}
\begin{itemize}
    \item \textbf{Definition:} Measures ability to discriminate correct vs incorrect predictions using uncertainty score
    \item \textbf{Range:} [0.5, 1.0] (0.5 = random, 1.0 = perfect)
    \item \textbf{Threshold:} $\geq$ 0.70 for viable selective prediction (set based on prior work)
    \item \textbf{Why:} Directly measures selective prediction quality---higher AUROC means better rejection of incorrect predictions
\end{itemize}

\textbf{Inference Cost (FLOPs-normalized):}
\begin{itemize}
    \item \textbf{Definition:} Total FLOPs per query, normalized to single forward pass = 1.0$\times$
    \item \textbf{Measurement:} MC dropout $k=N$ measured as $N.0\times \pm 0.1\times$
    \item \textbf{Why:} Hardware-agnostic cost metric (vs wall-clock time which varies with GPU/batching)
\end{itemize}

\textbf{Statistical Significance:}
\begin{itemize}
    \item \textbf{Method:} Paired $t$-test (two-tailed, $\alpha = 0.05$)
    \item \textbf{Application:} Pairwise AUROC comparisons, Pareto dominance testing
    \item \textbf{Sample size:} $n=3$ seeds (conservative power)
\end{itemize}

\textbf{Spearman Correlation ($\rho$):}
\begin{itemize}
    \item \textbf{Secondary metric:} Measures monotonic relationship between uncertainty score and incorrectness
    \item \textbf{Threshold:} $\rho > 0.2$ (validation check from h-m-integrated prerequisite)
    \item \textbf{Purpose:} Confirms uncertainty scores correlate with prediction correctness
\end{itemize}

All results reported as mean $\pm$ std across 3 random seeds.
